%% ── 14. TRABAJO FUTURO ───────────────────────────────────────────────────────
\chapter{Trabajo Futuro}

Este capítulo describe la hoja de ruta de evolución de ContractIA desde su
estado actual de MVP hasta una plataforma institucional completa. El trabajo
futuro se organiza en tres horizontes temporales articulados con el roadmap
del Capítulo 5, priorizando las acciones de mayor impacto y menor riesgo
técnico en cada etapa. Cada iniciativa incluye su motivación, entregables
esperados e impacto proyectado sobre los indicadores del sistema.

\section*{Corto Plazo (0--3 meses): Robustez y Cobertura Analítica}

El objetivo de este horizonte es consolidar el núcleo analítico y ampliar
su cobertura, partiendo de los resultados y limitaciones identificados en
la evaluación del MVP.

\begin{enumerate}

    \item \textbf{Agente Scout — RAG Agéntico (v10.0)}

    \textit{Motivación:} el RAG actual es estático: una única llamada de
    recuperación por sección, antes del análisis de agentes. Un Scout
    agéntico decide dinámicamente qué recuperar según lo que lee en cada
    fragmento del contrato, ejecutando múltiples rondas de recuperación
    si es necesario.

    \textit{Diseño:} loop ReAct manual (sin \texttt{bind\_tools} de
    LangChain por incompatibilidad con \texttt{with\_structured\_output});
    herramientas \texttt{buscar\_en\_contrato(consulta)} y
    \texttt{obtener\_clausula(numero)}; máximo 2 iteraciones por defecto;
    activado mediante flag \texttt{AGENTIC\_RAG\_ENABLED=true} para evitar
    regresión.

    \textit{Impacto esperado:} mejora del contexto entregado a los agentes
    en secciones con múltiples referencias cruzadas; reducción de falsos
    negativos en \texttt{OMISION\_NORMATIVA}.

    \item \textbf{Validación multi-contrato (Generalización)}

    \textit{Motivación:} los indicadores de precisión ($\approx 87\%$) y
    recall ($\approx 94\%$) fueron calculados sobre un único contrato piloto.
    La generalización del sistema requiere evidencia estadística sobre
    contratos de distinto sector, año y nivel de complejidad.

    \textit{Acción:} ejecutar el pipeline completo sobre 5--10 contratos de
    concesión de ProInversión de diferentes sectores (agua, energía, vial,
    telecomunicaciones) y comparar: número de hallazgos, distribución de
    tipos, tiempo de procesamiento y tasa de error de segmentación.

    \textit{Entregable:} informe de generalización con rangos de precisión
    y recall por sector; identificación de contratos problemáticos para el
    módulo anti-TOC.

    \item \textbf{Evaluación de modelos Claude (Anthropic via Vertex AI)}

    \textit{Motivación:} la solicitud de cuota para Claude Sonnet~4.6 y
    Opus~4.6 fue rechazada durante el periodo de evaluación. Los resultados
    preliminares en análisis semántico sugieren que los modelos Anthropic
    podrían ofrecer mayor precisión en la detección de inconsistencias
    lógicas sutiles.

    \textit{Acción:} reenviar solicitud de cuota a Vertex AI con
    justificación técnica detallada y caso de uso institucional; una vez
    aprobada, ejecutar benchmark completo sobre el Contrato A
    con Claude Sonnet~4.6 y Opus~4.6 bajo el mismo protocolo de evaluación.

    \textit{Entregable:} tabla comparativa de 4 modelos (Gemini~2.5~Pro,
    Gemini~3.1~Preview, Claude~Sonnet~4.6, Claude~Opus~4.6) en hallazgos,
    precisión, recall, latencia y costo.

    \item \textbf{Incremento de concurrencia y pruebas de carga}

    \textit{Motivación:} el MVP limita a 1 análisis concurrente
    (\texttt{Semaphore(1)}). En un escenario institucional con 3--5 auditores
    activos simultáneamente, la espera máxima podría superar 3 horas.

    \textit{Acción:} incrementar el semáforo a 3; ejecutar pruebas de carga
    con 3 contratos simultáneos en Cloud Run; medir impacto en cuota
    de Vertex AI, latencia de Cloud SQL y uso de memoria del contenedor.

    \textit{Entregable:} configuración validada con Semaphore(3); guía de
    escalado para la Fase 3 del roadmap.

    \item \textbf{Ampliación de la base normativa}

    \textit{Motivación:} la base actual cubre solo D.Leg.\ 1362, su
    Reglamento y normativa SUNASS. Los contratos de otros sectores
    referencian normativa no cubierta (MEM, MTC, Osiptel), lo que genera
    falsos negativos en \texttt{OMISION\_NORMATIVA}.

    \textit{Acción:} incorporar normativa sectorial de energía (MEM),
    transporte (MTC) y telecomunicaciones (Osiptel); desarrollar un
    pipeline automatizado de actualización periódica de la base vectorizada
    que detecte y cargue nuevas resoluciones y decretos.

\end{enumerate}

\section*{Mediano Plazo (3--6 meses): Plataforma Multiusuario y Análisis Avanzado}

El objetivo de este horizonte es transformar el MVP en una plataforma
accesible para múltiples perfiles institucionales, con capacidades de
análisis longitudinal y comparativo.

\begin{enumerate}

    \item \textbf{Arquitectura de validador dual (Cobertura + Precisión)}

    \textit{Motivación:} la paradoja cobertura--precisión observada entre
    Gemini~2.5~Pro y Gemini~3.1~Preview sugiere una arquitectura óptima:
    usar el modelo de mayor cobertura para el screening inicial y el modelo
    de mayor precisión para validar los hallazgos ALTA antes de incluirlos
    en el reporte final.

    \textit{Diseño:}
    \begin{itemize}
        \item Fase 1 — Gemini~2.5~Pro: análisis completo, todos los
        hallazgos (alta sensibilidad).
        \item Fase 2 — Gemini~3.1~Preview (o Claude~Opus~4.6): validación
        de los top-N hallazgos ALTA por sección (alta precisión).
        \item Hallazgos confirmados por ambas fases → severidad ALTA
        confirmada; resto → severidad MEDIA para revisión manual.
    \end{itemize}

    \textit{Impacto esperado:} reducción de falsos positivos del $\approx 13\%$
    al $< 5\%$ en hallazgos ALTA; mayor confianza institucional en el reporte.

    \item \textbf{Sistema de roles granular y auditoría de acciones}

    \textit{Motivación:} el MVP tiene dos roles básicos
    (\texttt{basico}/\texttt{auditor}). Una plataforma institucional
    requiere roles diferenciados con permisos específicos.

    \textit{Diseño:} roles propuestos: \texttt{legal} (acceso a hallazgos
    normativos), \texttt{tecnico} (acceso a hallazgos de plazos y lógica),
    \texttt{ambiental} (acceso a cláusulas ambientales), \texttt{auditor}
    (acceso completo), \texttt{director} (solo reporte ejecutivo),
    \texttt{admin} (gestión de usuarios y base normativa).
    Log inmutable de acciones de usuario en tabla \texttt{audit\_actions}.

    \item \textbf{Módulo de comparación entre versiones del contrato}

    \textit{Motivación:} los contratos de concesión son documentos vivos
    que se modifican mediante adendas. Detectar qué inconsistencias son
    nuevas, cuáles persisten y cuáles fueron corregidas entre versiones
    es una necesidad operativa crítica para la supervisión contractual.

    \textit{Funcionalidad:}
    \begin{itemize}
        \item Diff semántico entre dos versiones del mismo contrato:
        hallazgos nuevos, resueltos y persistentes.
        \item Vista ``track changes'' de obligaciones modificadas.
        \item Alerta automática si una adenda introduce una nueva
        inconsistencia en secciones previamente limpias.
    \end{itemize}

    \item \textbf{Panel de riesgos longitudinal y scoring contractual}

    \textit{Motivación:} los directivos de ProInversión necesitan una
    vista de riesgo agregada por proyecto, no hallazgo a hallazgo.

    \textit{Funcionalidad:}
    \begin{itemize}
        \item \textbf{Score de riesgo contractual} (0--100): ponderación
        de hallazgos por severidad, tipo y sección de impacto.
        \item \textbf{Dashboard por proyecto}: evolución del score entre
        versiones del contrato; comparación entre proyectos del mismo
        sector.
        \item \textbf{Top-10 hallazgos}: los de mayor impacto financiero
        o legal, priorizados automáticamente para revisión directiva.
    \end{itemize}

    \item \textbf{Integración con repositorio documental de ProInversión}

    \textit{Motivación:} el MVP requiere carga manual de contratos. Una
    integración con el repositorio documental institucional permitiría
    el análisis automático de nuevos contratos al momento de su publicación.

    \textit{Diseño:} endpoint \texttt{/webhooks/nuevo-contrato} que recibe
    notificación del repositorio; descarga automática del PDF; encolamiento
    inmediato del análisis; devolución del reporte al sistema de origen.

\end{enumerate}

\section*{Largo Plazo (6--12 meses): Escalabilidad, Certificación y Expansión}

El objetivo de este horizonte es convertir ContractIA en un estándar
institucional de auditoría contractual, certificado, escalable y extensible
a nuevos dominios.

\begin{enumerate}

    \item \textbf{Contenerización y orquestación para producción a escala}

    \textit{Motivación:} Cloud Run con \texttt{Semaphore(3)} es suficiente
    para el piloto institucional, pero no para una cartera de 50+ contratos
    concurrentes.

    \textit{Acción:} migración a Google Kubernetes Engine (GKE) o
    ECS-Fargate; procesamiento paralelo de contratos en colas distribuidas
    (Cloud Tasks o Pub/Sub); caché distribuido de embeddings normativos
    (Redis); autoescalado horizontal basado en la longitud de la cola.

    \item \textbf{FinOps IA: optimización de costos operativos}

    \textit{Motivación:} a escala de 100+ análisis/año, el costo de tokens
    LLM puede superar los \$10{,}000 anuales. La optimización de costos
    es un requisito para la sostenibilidad financiera del servicio.

    \textit{Acciones:}
    \begin{itemize}
        \item Caché de embeddings normativos: vectorizar la base normativa
        una sola vez y reutilizarla entre análisis.
        \item Reducción adaptativa de contexto: ajustar el tamaño del
        fragmento enviado al LLM según la complejidad estimada de la
        sección.
        \item Evaluación de modelos de menor costo para secciones de baja
        complejidad (Gemini Flash o Haiku como pre-filtro).
        \item Batch processing nocturno para análisis de bajo impacto temporal.
    \end{itemize}

    \item \textbf{Modelos predictivos de riesgo contractual}

    \textit{Motivación:} con un historial de 50+ contratos analizados,
    es posible construir modelos que anticipen el tipo y la densidad de
    inconsistencias esperadas antes de leer el contrato completo.

    \textit{Diseño:} features: sector, año, concedente, tipo de APP,
    extensión, número de secciones, historial de adendas. Target:
    score de riesgo y tipos de inconsistencia más probables. Uso:
    priorización automática de contratos para revisión urgente.

    \item \textbf{Certificación externa y validación metodológica}

    \textit{Motivación:} para que ContractIA sea adoptado como herramienta
    oficial en ProInversión, requiere validación por auditores o expertos
    legales independientes que certifiquen la metodología y los resultados.

    \textit{Acción:} comisionar una auditoría técnica externa sobre el
    pipeline de análisis; validación de precisión y recall sobre un conjunto
    de contratos con \textit{gold standard} anotado por especialistas;
    publicación de informe de certificación como aval institucional.

    \item \textbf{Extensión a nuevos dominios contractuales}

    \textit{Motivación:} la arquitectura multi-agente con RAG híbrido es
    generalizable a cualquier tipo de contrato de alta complejidad.

    \textit{Dominios priorizados:}
    \begin{itemize}
        \item \textbf{Contratos de obra pública (OSCE):} bajo Ley de
        Contrataciones del Estado; alta densidad de plazos y penalidades.
        \item \textbf{Concesiones mineras:} contratos de estabilidad
        jurídica y convenios de inversión con el Estado peruano.
        \item \textbf{Contratos de suministro de energía (COES/OSINERGMIN):}
        formulación financiera compleja; alto impacto regulatorio.
        \item \textbf{Contratos bilaterales de inversión (BIT):}
        documentos en inglés o español técnico-diplomático.
    \end{itemize}

    \textit{Adaptaciones requeridas:} prompts específicos por tipo de
    contrato; base normativa sectorial; esquemas de clasificación de
    hallazgos adaptados.

\end{enumerate}

\section*{Visión a 18+ meses: ContractIA como Plataforma SaaS}

En el horizonte de 18 meses y más, ContractIA tiene el potencial de
evolucionar desde una herramienta institucional interna hacia una
\textbf{plataforma SaaS de auditoría contractual} orientada a entidades
públicas y privadas de la región:

\begin{itemize}
    \item \textbf{Modelo de negocio:} suscripción por volumen de análisis
    (contratos/mes); acceso API para integración con sistemas de terceros.
    \item \textbf{Expansión regional:} Colombia (ANI), Chile (MOP), México
    (SCT), con base normativa y prompts adaptados a cada marco legal nacional.
    \item \textbf{Marketplace de módulos normativos:} organizaciones que
    generan sus propias bases normativas y las comparten en el ecosistema.
    \item \textbf{Integración con firmas legales y consultoras:} como
    herramienta de due diligence acelerada en fusiones y adquisiciones.
\end{itemize}

\section*{Resumen del Roadmap}

\begin{table}[H]
\centering
\caption{Resumen del roadmap de trabajo futuro}
\label{tab:roadmap}
\begin{tabular}{p{1.5cm} p{3cm} p{4cm} p{4.5cm}}
\toprule
\textbf{Plazo} & \textbf{Foco} & \textbf{Iniciativas clave} & \textbf{Impacto esperado} \\
\midrule
0--3 m &
    Robustez y cobertura &
    Scout agéntico; multi-contrato; Claude; concurrencia$\times$3; base normativa &
    Precisión $+5\%$; recall $+3\%$; evaluación de 4 modelos LLM \\
3--6 m &
    Plataforma multiusuario &
    Validador dual; roles; diff versiones; scoring; integración documental &
    FP ALTA $< 5\%$; dashboard directivo; análisis automático \\
6--12 m &
    Escalabilidad y certificación &
    Kubernetes; FinOps IA; modelos predictivos; certificación; nuevos dominios &
    100+ análisis/año; costo $-40\%$; aval institucional \\
18+ m &
    Plataforma regional &
    SaaS; expansión regional; marketplace normativo &
    Monetización; impacto en 3+ países \\
\bottomrule
\end{tabular}
\end{table}